\section{Resultados e discussão}
\label{sec:resultados}

\subsection{Comparação geral}

A Tabela~\ref{tab:resumo} reúne a configuração e os principais indicadores de todas as execuções. A Tabela~\ref{tab:perdas} detalha os denominadores por tarefa, evitando que a taxa hard agregada oculte uma perda de segurança. Os tempos completos, incluindo mínimos, médias, máximos, jitter e CPU por tarefa, constam do Apêndice~\ref{ap:metricas}.

\begin{singlespace}
\begingroup
\fontsize{8}{10}\selectfont
\setlength{\tabcolsep}{3pt}
\renewcommand{\arraystretch}{1.18}
\begin{longtable}{@{}lllrrrrrrrr@{}}
\caption{Comparação dos 18 cenários: configuração, perdas, segurança e CPU.}\label{tab:resumo}\\
\toprule
ID & Política & Modo & MHz & Slice & s & H (\%) & S (\%) & $L_{FS}$ & $R_{FS}$ & CPU (\%)\\
\midrule
\endfirsthead
\multicolumn{11}{c}{\tablename\ \thetable\ -- continuação}\\
\toprule
ID & Política & Modo & MHz & Slice & s & H (\%) & S (\%) & $L_{FS}$ & $R_{FS}$ & CPU (\%)\\
\midrule
\endhead
\midrule\multicolumn{11}{r}{Continua na próxima página.}\\
\endfoot
\bottomrule
\endlastfoot
01 & CUSTOM & PRE & 80 & ON & 35{,}000 & 0{,}007 & 41{,}667 & 0{,}020 & 4{,}259 & 69{,}863\\
02 & CUSTOM & PRE & 240 & ON & 35{,}000 & 0{,}000 & 0{,}000 & 0{,}009 & 1{,}348 & 23{,}210\\
03 & CUSTOM & COOP & 80 & ON & 35{,}015 & 0{,}208 & 41{,}667 & 0{,}021 & 4{,}799 & 69{,}811\\
04 & CUSTOM & COOP & 240 & ON & 35{,}012 & 0{,}136 & 0{,}000 & 0{,}008 & 1{,}347 & 23{,}183\\
05 & DM & PRE & 80 & ON & 35{,}000 & 0{,}007 & 41{,}667 & 0{,}024 & 14{,}792 & 69{,}816\\
06 & DM & PRE & 240 & ON & 35{,}000 & 0{,}000 & 0{,}000 & 0{,}013 & 3{,}162 & 23{,}195\\
07 & DM & COOP & 80 & ON & 35{,}015 & 0{,}521 & 41{,}667 & 1{,}420 & 9{,}042 & 69{,}976\\
08 & DM & COOP & 240 & ON & 35{,}012 & 0{,}114 & 0{,}000 & 0{,}013 & 1{,}351 & 23{,}182\\
E01 & CUSTOM & PRE & 160 & ON & 35{,}000 & 0{,}000 & 0{,}000 & 0{,}012 & 2{,}022 & 34{,}720\\
E02 & CUSTOM & COOP & 160 & ON & 35{,}013 & 0{,}179 & 0{,}000 & 0{,}024 & 2{,}047 & 34{,}708\\
E03 & DM & PRE & 160 & ON & 34{,}999 & 0{,}000 & 0{,}000 & 0{,}033 & 3{,}781 & 34{,}722\\
E04 & DM & COOP & 160 & ON & 35{,}013 & 0{,}186 & 0{,}000 & 0{,}017 & 2{,}026 & 34{,}761\\
E05 & DM & PRE & 80 & OFF & 35{,}003 & 0{,}007 & 41{,}667 & 3{,}517 & 18{,}983 & 69{,}795\\
E06 & DM & PRE & 240 & OFF & 30{,}001 & 0{,}000 & 0{,}000 & 0{,}013 & 1{,}351 & 23{,}262\\
E07 & RM & PRE & 240 & ON & 30{,}000 & 0{,}000 & 0{,}000 & 0{,}530 & 1{,}905 & 23{,}246\\
E08 & RM & PRE & 80 & ON & 30{,}000 & 0{,}008 & 41{,}667 & 0{,}948 & 16{,}871 & 69{,}971\\
E09 & RM & COOP & 80 & ON & 30{,}015 & 0{,}592 & 41{,}667 & 0{,}024 & 8{,}042 & 70{,}139\\
E10 & RM & COOP & 240 & ON & 35{,}012 & 0{,}114 & 0{,}000 & 0{,}025 & 1{,}367 & 23{,}192\\
\end{longtable}
\endgroup
\noindent{\footnotesize Fonte: logs preservados; processamento do autor. PRE: preemptivo; COOP: configuração cooperativa. Todos UNICORE. Slicing declarado conforme histórico, não impresso no banner. Hard: FUS+CTRL+FS; soft: NAV (B+C). L e R de FS em ms; R termina após a carga artificial. CPU é a média aproximada da soma das quatro tarefas, não CPU total.}
\end{singlespace}

\begin{singlespace}
\begingroup
\footnotesize
\setlength{\tabcolsep}{3pt}
\renewcommand{\arraystretch}{1.18}
\begin{longtable}{@{}lrrrr@{}}
\caption{Perdas por tarefa: misses/jobs e percentual acumulado.}\label{tab:perdas}\\
\toprule
ID & FUS & CTRL & NAV & FS\\
\midrule
\endfirsthead
\multicolumn{5}{c}{\tablename\ \thetable\ -- continuação}\\
\toprule
ID & FUS & CTRL & NAV & FS\\
\midrule
\endhead
\midrule\multicolumn{5}{r}{Continua na próxima página.}\\
\endfoot
\bottomrule
\endlastfoot
01 & 0/7000 (0{,}000\%) & 1/6999 (0{,}014\%) & 5/12 (41{,}667\%) & 0/1 (0{,}000\%)\\
02 & 0/7000 (0{,}000\%) & 0/7000 (0{,}000\%) & 0/12 (0{,}000\%) & 0/1 (0{,}000\%)\\
03 & 26/7001 (0{,}371\%) & 3/6967 (0{,}043\%) & 5/12 (41{,}667\%) & 0/1 (0{,}000\%)\\
04 & 19/7001 (0{,}271\%) & 0/6974 (0{,}000\%) & 0/12 (0{,}000\%) & 0/1 (0{,}000\%)\\
05 & 0/7000 (0{,}000\%) & 0/7000 (0{,}000\%) & 5/12 (41{,}667\%) & 1/1 (100{,}000\%)\\
06 & 0/7000 (0{,}000\%) & 0/7000 (0{,}000\%) & 0/12 (0{,}000\%) & 0/1 (0{,}000\%)\\
07 & 48/7001 (0{,}686\%) & 25/7001 (0{,}357\%) & 5/12 (41{,}667\%) & 0/1 (0{,}000\%)\\
08 & 16/7001 (0{,}229\%) & 0/7001 (0{,}000\%) & 0/12 (0{,}000\%) & 0/1 (0{,}000\%)\\
E01 & 0/7000 (0{,}000\%) & 0/7000 (0{,}000\%) & 0/12 (0{,}000\%) & 0/1 (0{,}000\%)\\
E02 & 25/7001 (0{,}357\%) & 0/6975 (0{,}000\%) & 0/12 (0{,}000\%) & 0/1 (0{,}000\%)\\
E03 & 0/7000 (0{,}000\%) & 0/7000 (0{,}000\%) & 0/12 (0{,}000\%) & 0/1 (0{,}000\%)\\
E04 & 26/7001 (0{,}371\%) & 0/7001 (0{,}000\%) & 0/12 (0{,}000\%) & 0/1 (0{,}000\%)\\
E05 & 0/7000 (0{,}000\%) & 0/7000 (0{,}000\%) & 5/12 (41{,}667\%) & 1/1 (100{,}000\%)\\
E06 & 0/6000 (0{,}000\%) & 0/6000 (0{,}000\%) & 0/12 (0{,}000\%) & 0/1 (0{,}000\%)\\
E07 & 0/6000 (0{,}000\%) & 0/6000 (0{,}000\%) & 0/12 (0{,}000\%) & 0/1 (0{,}000\%)\\
E08 & 0/6000 (0{,}000\%) & 0/6000 (0{,}000\%) & 5/12 (41{,}667\%) & 1/1 (100{,}000\%)\\
E09 & 47/6001 (0{,}783\%) & 24/6001 (0{,}400\%) & 5/12 (41{,}667\%) & 0/1 (0{,}000\%)\\
E10 & 16/7001 (0{,}229\%) & 0/7001 (0{,}000\%) & 0/12 (0{,}000\%) & 0/1 (0{,}000\%)\\
\end{longtable}
\endgroup
\noindent{\footnotesize Fonte: logs preservados. Percentuais reportados pelo firmware. O contador da FUS depende da referência discutida na Seção~\ref{sec:limitacoes}.}
\end{singlespace}


Seis execuções não registraram perdas nas quatro tarefas: CUSTOM preemptivo a 160/240~MHz, DM preemptivo a 160/240~MHz, DM preemptivo a 240~MHz com slicing OFF e RM preemptivo a 240~MHz. Em todas as condições avaliadas a 80~MHz, NAV registrou 5/12 perdas (41,667\%). Todos os ensaios configurados como cooperativos registraram perdas da FUS. A ausência de perdas em uma execução não demonstra uma garantia hard RT para todas as fases de entrada.

\figura{01_perdas_principal.pdf}{Perdas de deadline por tarefa na matriz principal.}{fig:perdas}{As células apresentam misses/jobs e percentual. As cores indicam ocorrência de perdas, sem equiparar sua gravidade entre tarefas. FUS utiliza a referência temporal discutida na Seção~\ref{sec:limitacoes}.}

\subsection{Prioridades: CUSTOM e DM}

A 80~MHz no modo preemptivo, CUSTOM registrou uma perda do CTRL e nenhuma de FS; DM registrou zero perdas do CTRL e uma de FS. As respostas da emergência foram 4,259~ms e 14,792~ms, respectivamente, para D=10~ms. Esse comportamento é compatível com a prioridade máxima de FS em CUSTOM e com FUS/CTRL acima de FS em DM.

No CTRL de CUSTOM/80, a duração média foi 1,3561~ms, enquanto a máxima chegou a 10,165~ms e a resposta máxima a 12,223~ms. Em DM/80, a resposta máxima do CTRL foi 3,795~ms. Portanto, a média reduzida não elimina um atraso pontual importante. Sem trace de trocas de contexto, não é possível atribuir cada intervalo dessa duração a uma interferência específica.

A 240~MHz, CUSTOM e DM preemptivos registraram zero perdas em todas as tarefas. A resposta de FS foi 1,348~ms em CUSTOM e 3,162~ms em DM. A maior capacidade de processamento proporcionou folga observada para os deadlines nessa sequência. A Figura~\ref{fig:respostas} apresenta os máximos e seus limites.

\figura{02_respostas_deadlines.pdf}{Respostas máximas e deadlines na matriz principal.}{fig:respostas}{As linhas tracejadas indicam os deadlines configurados. Os máximos são observados; em FS o fim é posterior à escrita dos motores zerados.}

\subsection{Resposta da emergência}

A Figura~\ref{fig:fs} separa a latência até iniciar da duração do trecho. Como há um único job FS por execução, as duas parcelas pertencem ao mesmo evento e podem ser somadas. Em CUSTOM preemptivo, L foi de 20, 12 e 9~$\mu$s para 80, 160 e 240~MHz, respectivamente; R foi de 4,259, 2,022 e 1,348~ms.

Em DM preemptivo/80, FS iniciou após 24~$\mu$s, mas terminou 14,792~ms após o evento. A duração decorrida foi 14,768~ms. Assim, o miss não significa uma espera de quase 15~ms para começar: a tarefa começou cedo e o trecho completo terminou depois do prazo. Como FUS e CTRL têm maior prioridade, interferência durante esse trecho é uma explicação compatível com a implementação.

Os motores são escritos como zero antes da carga artificial, e esse instante não tem timestamp próprio. O resultado demonstra atraso de conclusão do job instrumentado, mas não permite afirmar que a ação de zerar os motores ocorreu somente ao final. Também não mede a latência anterior à detecção do touch. A observação 1/1 de FS descreve o evento coletado; não estima uma probabilidade de falha em operação prolongada.

\figura{03_fail_safe_todos.pdf}{Decomposição da resposta de segurança nos 18 ensaios.}{fig:fs}{L é o intervalo ISR--início; E é o intervalo início--fim instrumentado. A referência de 10~ms aplica-se ao deadline contado pelo firmware. Cada barra representa um evento de emergência.}

\subsection{Preempção, cooperação e jitter}

Na matriz principal, CUSTOM cooperativo apresentou 26/7001 perdas de FUS a 80~MHz e 19/7001 a 240~MHz. DM cooperativo apresentou 48/7001 e 16/7001. O CTRL perdeu 3/6967 e 25/7001 jobs a 80~MHz, em CUSTOM e DM, respectivamente; a 240~MHz não registrou perdas.

Trechos longos sem preempção automática podem atrasar tarefas prontas até um ponto de bloqueio ou yield. Aumentar a frequência reduz o custo desses trechos, mas não substitui a análise dos pontos de reescalonamento. Além disso, a referência da FUS mistura ticks e microssegundos; por isso seus valores são apresentados como contadores reportados, com a ressalva da Seção~\ref{sec:limitacoes}.

O FS de DM cooperativo/80 concluiu em 9,042~ms, abaixo dos 14,792~ms do preemptivo. Essa observação não demonstra superioridade geral do cooperativo: a fase do evento varia e o caminho da ISR pode solicitar uma troca de tarefa. O trecho após o início também pode sofrer menos interrupções de tarefas superiores nessa configuração.

A Figura~\ref{fig:jitter} mostra desvios máximos de início da FUS próximos de um período em vários cenários cooperativos. Os painéis distinguem a definição de jitter da FUS da variação de latência entre jobs do CTRL. A Figura~\ref{fig:evolucao} apresenta a taxa de perdas por janela, obtida por diferenças de contadores acumulados. Ela localiza os aumentos em intervalos de aproximadamente 5~s, sem reconstruir o instante exato de cada perda ou isolar causalmente o efeito do touch A.

\figura{06_jitter.pdf}{Jitter máximo reportado na matriz principal.}{fig:jitter}{FUS: desvio absoluto da referência esperada, sujeito ao desalinhamento de fase. CTRL: diferença absoluta entre latências de jobs consecutivos. As duas estatísticas têm definições distintas.}
\figura{07_evolucao_perdas.pdf}{Evolução das perdas da FUS nas configurações cooperativas.}{fig:evolucao}{Cada série corresponde a uma execução. As taxas são por janela, não médias de repetições; fases dos eventos manuais diferem entre cenários.}

\subsection{Frequência e extra de 160 MHz}

Para CUSTOM preemptivo, o mínimo observado de E da FUS foi 2,027~ms a 80~MHz, 1,008~ms a 160~MHz e 0,668~ms a 240~MHz. No CTRL, os mínimos foram 1,352, 0,670 e 0,447~ms. A redução é compatível com um número fixo de cálculos executado em frequência maior. Esses mínimos caracterizam o custo base observado e não WCET.

O extra de 160~MHz acrescentou quatro execuções. CUSTOM e DM preemptivos não registraram perdas; os cooperativos registraram 25 e 26 perdas de FUS, respectivamente, sem perdas de CTRL, NAV ou FS. Assim, a frequência intermediária já eliminou, nas fases observadas, as perdas de NAV e dos preemptivos. Isso não estabelece 160~MHz como frequência mínima universalmente suficiente.

\figura{04_frequencia.pdf}{Efeito observado da frequência no custo base e na emergência.}{fig:frequencia}{FUS/CTRL usam duração mínima observada; FS usa resposta do único evento. Linhas conectam medições, sem representar um modelo ajustado ou previsão para frequências não testadas.}

A soma média aproximada da CPU das quatro tarefas ficou em torno de 70\% a 80~MHz, 35\% a 160~MHz e 23\% a 240~MHz (Figura~\ref{fig:cpu}). FUS e CTRL dominam o tempo acumulado devido à taxa nominal de 200~Hz. NAV tem baixa ocupação média, mas demanda concentrada em poucos eventos; sua pequena média não impede uma resposta máxima longa.

Na matriz principal, a resposta máxima de NAV variou de 46,686 a 69,521~ms a 80~MHz, acima de D=20~ms; a 240~MHz variou de 7,424 a 8,672~ms. O caminho de rota possui carga longa e prioridade inferior às tarefas críticas. A taxa 5/12 é da tarefa que reúne B e C; não pode ser automaticamente apresentada como taxa exclusiva de comandos de rota. O enunciado não fixa percentual tolerável para soft RT, portanto a taxa encontrada deve ser descrita como degradação de serviço, sem declará-la aceitável por definição.

\figura{05_cpu_media.pdf}{Ocupação média aproximada atribuída às quatro tarefas.}{fig:cpu}{Média ponderada pelas janelas impressas. A soma exclui categorias próprias para tarefas auxiliares e sistema; seu complemento não foi interpretado como IDLE.}

Uma ocupação média abaixo de 100\% não garante deadlines curtos sob demandas concentradas, bloqueios e prioridades fixas. Também não foi medida energia elétrica; os resultados de frequência e CPU não quantificam economia de energia.

\subsection{Extra de time slicing}

Os cenários E05/E06 repetem DM preemptivo a 80/240~MHz com slicing declarado OFF. FUS e CTRL têm a mesma prioridade, condição em que o rodízio por tick pode mudar sua alternância. Desligar slicing não elimina bloqueios, yields ou preempção por prioridades superiores~\cite{barry_freertos}.

A 80~MHz, a resposta de FS passou de 14,792~ms com ON para 18,983~ms com OFF; L passou de 0,024 para 3,517~ms. Ambos registraram um miss de FS e cinco de NAV, com zero perdas de FUS/CTRL. A 240~MHz, R de FS foi 3,162~ms com ON e 1,351~ms com OFF, sem perdas em ambas as execuções.

O sentido diferente dessa variação entre frequências impede concluir que desligar slicing é sempre vantajoso ou prejudicial. Houve um evento por condição, sem sincronização de fase. O valor de slicing não foi impresso no banner e depende do histórico da coleta, limitando a confirmação retrospectiva.

\figura{08_time_slicing.pdf}{Comparação exploratória de time slicing em DM preemptivo.}{fig:slicing}{Os valores ON/OFF são declarados no histórico. As diferenças observadas não isolam o efeito causal do rodízio, devido à ausência de repetições e de controle de fase.}

\subsection{Extra RM}

RM preemptivo/240 registrou zero perdas; a 80~MHz, registrou 5/12 em NAV e 1/1 em FS, cuja resposta foi 16,871~ms. RM cooperativo/80 apresentou 47/6001 perdas de FUS e 24/6001 de CTRL; a 240~MHz, 16/7001 de FUS e zero nas demais tarefas. A proximidade com DM é compatível com a tabela de prioridades idêntica (Figura~\ref{fig:rm}).

As durações diferentes tornam inadequada a comparação de contagens isoladas: DM cooperativo/80 teve 48 perdas de FUS em 7001 jobs, enquanto RM teve 47 em 6001. As taxas finais são 0,686\% e 0,783\%. Em resumos próximos de 30~s, os contadores foram 45/6001 e 47/6001, respectivamente. No cooperativo/240, DM e RM apresentaram 14/6001 nesse recorte. Não há evidência suficiente para atribuir superioridade a uma política, sobretudo porque sua ordenação implementada coincide.

\figura{09_dm_rm.pdf}{DM e RM: comparação das taxas reportadas.}{fig:rm}{As prioridades implementadas são iguais. As taxas usam o denominador de cada execução; durações e fases dos eventos diferem. O tratamento agregado de B+C limita uma interpretação clássica de RM.}

\subsection{Defasagem entre fusão e controle}

Os snapshots finais de CUSTOM cooperativo apresentam diferenças de 34, 26 e 27 jobs entre FUS e CTRL, a 80, 160 e 240~MHz (Figura~\ref{fig:fila}). A fila de uma posição mantém a amostra mais recente, podendo substituir estados ainda não consumidos. A diferença também pode incluir trabalho em trânsito no instante do resumo; não é um contador exato de amostras perdidas.

Os misses de CTRL são contados sobre jobs concluídos. Portanto, zero misses não demonstra consumo de todas as amostras. Seria necessário registrar os saltos de sequência e o trabalho pendente para separar as parcelas. O campo de descartes touch pertence a outro caminho de eventos e não cobre essa fila.

\figura{10_defasagem_fila.pdf}{Diferença de jobs entre FUS e CTRL no snapshot final.}{fig:fila}{Indicador de defasagem, que pode incluir sobrescrita e trabalho em trânsito. Não representa uma contagem exata de amostras descartadas.}

\FloatBarrier
